\documentclass[12pt,a4paper]{article}

% ============================================================
% PACKAGES
% ============================================================
\usepackage[
    a4paper,
    top=1.55cm,
    bottom=1.35cm,
    left=1.75cm,
    right=1.75cm
]{geometry}
\usepackage[T1]{fontenc}
\usepackage{newtxtext}
\usepackage{graphicx}
\usepackage{xcolor}
\usepackage{array}
\usepackage{tabularx}
\usepackage{fancyhdr}
\usepackage{tikz}
\usepackage{eso-pic}
\usepackage{ragged2e}
\usepackage{enumitem}

% ============================================================
% COURSE
% ============================================================
\newcommand{\CourseHeader}{OSDBMS (TCS-502 / TCS-503)}

% ============================================================
% COLORS
% ============================================================
\definecolor{GEUBlue}{RGB}{43,91,145}
\definecolor{GEUBannerBlue}{RGB}{0,102,148}
\definecolor{GEUYellow}{RGB}{255,201,38}
\definecolor{InstructionGrey}{RGB}{125,125,125}

% ============================================================
% GENERAL SETTINGS
% ============================================================
\setlength{\parindent}{0pt}
\setlength{\parskip}{0pt}
\renewcommand{\arraystretch}{1}
\setlist{leftmargin=1.3em,itemsep=1.5pt,topsep=3pt,parsep=0pt,partopsep=0pt}

% ============================================================
% HEADER / FOOTER
% ============================================================
\pagestyle{fancy}
\fancyhf{}
\fancyhead[L]{%
    \fontsize{10.5}{12}\selectfont
    \CourseHeader
}
\fancyhead[R]{%
    \fontsize{10.5}{12}\selectfont
    PHASE-III PROJECT PROGRESS
}
\fancyfoot[R]{%
    \fontsize{10}{12}\selectfont
    Page \thepage
}
\renewcommand{\headrulewidth}{0pt}
\renewcommand{\footrulewidth}{0pt}
\setlength{\headheight}{20pt}
\setlength{\headsep}{0.65cm}
\setlength{\footskip}{0.65cm}

% ============================================================
% BOTTOM YELLOW LINE
% ============================================================
\AddToShipoutPictureBG{%
\begin{tikzpicture}[remember picture,overlay]
    \fill[GEUYellow]
    (current page.south west)
    rectangle
    ([yshift=0.14cm]current page.south east);
\end{tikzpicture}
}

% ============================================================
% CUSTOM COMMANDS
% ============================================================
\newcommand{\sectiontitle}[1]{%
    {\color{GEUBlue}
    \fontsize{18}{22}\selectfont #1}
}
\newcommand{\instruction}[1]{%
    {\color{InstructionGrey}
    \fontsize{10.5}{13}\selectfont #1}
}
% The box grows with its text, so nothing spills outside the border.
% The first argument (the template's fixed height) is kept for compatibility.
\newcommand{\answerbox}[2]{%
    \noindent
    \begin{tabularx}{\textwidth}{
        |>{\raggedright\arraybackslash}X|}
        \hline
        \parbox[t]{\dimexpr\linewidth-12pt\relax}{
            \vspace{5pt}
            \fontsize{10.5}{13}\selectfont #2
            \par\vspace{6pt}
        }\\
        \hline
    \end{tabularx}
}
% A table cell with a little padding above and below.
\newcommand{\cell}[1]{\parbox[t]{\linewidth}{\vspace{3pt}\fontsize{10}{12.5}\selectfont #1\par\vspace{4pt}}}

% ============================================================
% DOCUMENT
% ============================================================
\begin{document}

% ============================================================
% PAGE 1
% ============================================================
\vspace*{0.10cm}
\includegraphics[width=3.2cm]{geulogo.png}

\vspace{0.40cm}
\noindent
\colorbox{GEUBannerBlue}{%
\parbox{\dimexpr\textwidth-2\fboxsep\relax}{
    \vspace{0.17cm}
    {\color{GEUYellow}
    \fontsize{16.5}{20}\selectfont
    \bfseries
    PROJECT AND TEAM INFORMATION}
    \vspace{0.17cm}
}}

\vspace{0.25cm}

% PROJECT TITLE
\sectiontitle{Project Title}

\vspace{0.30cm}
\answerbox{1.00cm}{OrchestrOS: Intelligent Container Orchestration with Resource-Aware Scheduling and
Autoscaling}

\vspace{0.45cm}

% ============================================================
% TEAM TABLE
% Each row: who the member is (left), and their details and photo (right).
% Responsibilities list the Phase-II work, then what was added in Phase III.
% ============================================================
\newcommand{\memberleft}[3]{%
\begin{minipage}[t]{\linewidth}
\vspace{1.5mm}
\fontsize{11}{13.2}\selectfont
#1\\
(Last Name, name: student ID: email, picture):\\[3pt]
Role: #2\\
Responsibilities: #3
\vspace{2mm}
\end{minipage}}
\newcommand{\memberright}[4]{%
\begin{minipage}[t]{\linewidth}
\vspace{1.5mm}
\fontsize{11}{13.2}\selectfont
#1 -- #2\\
#3\\[3pt]
\includegraphics[width=2.2cm,height=1.6cm,keepaspectratio]{#4}
\vspace{2mm}
\end{minipage}}

\noindent
\begin{tabularx}{\textwidth}{|>{\RaggedRight\arraybackslash}p{0.53\textwidth}|>{\RaggedRight\arraybackslash}X|}
\hline
\begin{minipage}[t]{\linewidth}
\vspace{1.5mm}
\fontsize{11}{13.2}\selectfont
Team Name:\\
Team \#
\vspace{1.5mm}
\end{minipage}
&
\begin{minipage}[t]{\linewidth}
\vspace{1.5mm}
\fontsize{11}{13.2}\selectfont
OSDBMS-V-T244
\vspace{1.5mm}
\end{minipage}
\\
\hline
\memberleft{Team member 1 (Team Lead)}{Database \& Testing (DB/Testing)}{MySQL 8 Schema, ACID Row-Lock
Reservations, Monitoring Queries \& Automated Tests; Phase III: Reactive Autoscaler, Reactive vs ML
Comparison Module \& Integration Test Suite}
&
\memberright{Chamoli, Pranjal}{24021077 (Roll: 2025368)}{24021077@geu.ac.in}{image1.png}
\\
\hline
\memberleft{Team member 2}{Developer}{OS Scheduling Policies (FCFS, SJF, Priority, Round Robin) \& Worker
Placement Engine; Phase III: Round Robin Preemption \& Requeue and Scheduling-Policy Benchmark}
&
\memberright{Rawat, Maitri}{240221432 (Roll: 2026063)}{240221432@geu.ac.in}{image2.png}
\\
\hline
\memberleft{Team member 3}{Developer}{Core Orchestrator Pipeline, Express REST API, Docker Engine Socket
Execution \& Crash Recovery; Phase III: ML-Assisted Autoscaling (Linear-Regression Forecaster) \&
Job Cancellation/Retries}
&
\memberright{Joshi, Nikhil}{240211642 (Roll: 2025339)}{240211642@geu.ac.in}{image3.png}
\\
\hline
\memberleft{Team member 4}{Developer, React Dashboard UI \& Workload Generator}{React 19 Operator UI,
Seeded Workload Generator v2 \& Project Documentation; Phase III: Autoscaling \& Comparison Dashboard
Pages and Workload Generator Refinements}
&
\memberright{Kunwar, Vansh}{24021068 (Roll: 2025560)}{24021068@geu.ac.in}{image4.png}
\\
\hline
\end{tabularx}

\newpage

% ============================================================
% PHASE III: ABSTRACT AND ARCHITECTURE
% ============================================================
\vspace*{0.15cm}
\noindent
\colorbox{GEUBannerBlue}{%
\parbox{\dimexpr\textwidth-2\fboxsep\relax}{
    \vspace{0.18cm}
    {\color{GEUYellow}
    \fontsize{17}{21}\selectfont
    \bfseries
    PHASE-III PROJECT PROGRESS DESCRIPTION (35 pts)}
    \vspace{0.18cm}
}}

\vspace{0.70cm}

% ABSTRACT
\sectiontitle{Project Abstract (2 pts)}

\vspace{0.05cm}
\instruction{(Brief restatement of your project's main goal. Max 300 words).}

\vspace{0.12cm}
\answerbox{3.60cm}{%
OrchestrOS is a local, Kubernetes-inspired workload orchestrator that applies core Operating Systems
and DBMS concepts in one working system. Users generate seeded, reproducible batches of jobs; each
job carries a CPU and memory request, a priority and an arrival time.

Every job moves through an explicit pipeline. A CPU scheduler picks the next job (FCFS, SJF, Priority,
or Round Robin with preemption), a placement strategy chooses a logical worker node, and a MySQL
transaction reserves that worker's capacity using row-level locks, so concurrent jobs can never
over-allocate it. The job then runs as an isolated Docker container limited to exactly the reserved CPU
and memory. Recording the result (with a checksum of the work done) releases the capacity in the same
transaction.

On top of this pipeline, OrchestrOS scales the cluster automatically. A \textbf{reactive} autoscaler adds
or retires workers based on the queue and reserved utilisation it measures now. An
\textbf{ML-assisted} strategy fits a linear-regression model to recent measurements and forecasts the
load 60 seconds ahead, so it can add capacity before the load arrives. Two built-in experiments compare
the four scheduling policies, and Reactive versus ML-assisted autoscaling, on identical,
checksum-verified workloads. They report waiting time, turnaround, makespan and extra capacity used.

A React dashboard controls every stage and shows live cluster state, charts and a plain-language
verdict for each experiment. The project is complete, and its comparison results feed the team's
research write-up on ML-based autoscaling.}

\vspace{0.70cm}

% UPDATED APPROACH
\sectiontitle{Updated Project Approach and Architecture (2 pts)}

\vspace{0.05cm}
\instruction{(Describe your current approach, including system design,
communication protocols, libraries used, etc.\\ Max 300 words).}

\vspace{0.25cm}
\answerbox{4.75cm}{%
OrchestrOS is a three-tier system that runs on one machine:
\begin{itemize}
    \item \textbf{Frontend:} a React 19 single-page dashboard built with Vite 8. It talks to the backend
    only through REST/JSON over HTTP and polls for live state (every 2\,s while an experiment runs,
    5\,s when idle). Charts are hand-built SVG; no chart library is used.
    \item \textbf{Backend:} Node.js with Express 5 and Zod request validation, split into 13 modules
    (jobs, workers, workloads, scheduler, placement, resources, executions, orchestrator, autoscaler,
    monitoring, benchmarks, scaling-comparisons, health). Each module has routes, a service with the
    logic, and a repository for data access.
    \item \textbf{Database:} MySQL 8.4 (InnoDB) through the Prisma 6 ORM, with 5 versioned migrations.
    Reservation, release, cancellation, requeue and scaling run as transactions that lock job and
    worker rows (\texttt{SELECT \ldots{} FOR UPDATE}). Deadlock victims are retried automatically, and
    \texttt{CHECK} constraints reject invalid execution records.
    \item \textbf{Execution:} the backend calls the Docker Engine API directly over the local socket.
    Each job runs in a purpose-built workload-runner image with no network, a read-only filesystem,
    all capabilities dropped, and CPU/memory limits equal to its reservation.
    \item \textbf{Autoscaling:} one autoscaler evaluates the cluster every 5\,s with one active
    strategy. Reactive applies thresholds to the measured queue and utilisation; ML fits
    least-squares linear regression over a sliding window of samples and acts on the 60\,s forecast.
    Every decision is stored with its reason.
    \item \textbf{Experiments:} a registry allows one benchmark or comparison at a time. Each run starts
    from a reset cluster with the same seed; differences under 5\% are reported as ties.
    \item \textbf{Deployment:} Docker Compose starts MySQL, a migration/seed step, the backend and
    the frontend.
\end{itemize}}

\newpage

% ============================================================
% TASKS AND CHALLENGES
% ============================================================
\vspace*{0.20cm}

% TASKS COMPLETED
\sectiontitle{Tasks Completed (7 pts)}

\vspace{0.05cm}
\instruction{(Describe the main tasks that have been assigned and already completed. Max 250 words).}

\vspace{0.20cm}
\noindent
\begin{tabularx}{\textwidth}{
|>{\raggedright\arraybackslash}p{0.66\textwidth}
|>{\raggedright\arraybackslash}X|
}
\hline
\textbf{Task Completed} &
\textbf{Team Member}
\\
\hline
\cell{\textbf{Phase II:} MySQL 8 schema and Prisma migrations; ACID row-lock reservations that prevent
over-allocation; monitoring queries; automated tests; move from TypeScript/PostgreSQL to
JavaScript/MySQL.\\[2pt]
\textbf{Phase III:} reactive autoscaler; MySQL deadlock retry; merging the ML forecaster into the
autoscaler as a strategy; Reactive vs ML comparison module (3 new tables); integration test suite
(241 backend tests).}
& \cell{Chamoli, Pranjal\\(Team Lead, DB/Testing)}
\\
\hline
\cell{\textbf{Phase II:} FCFS, SJF, Priority and Round Robin scheduling; worker placement engine
(least-loaded, first-fit, resource-aware).\\[2pt]
\textbf{Phase III:} Round Robin preemption with requeue; placement tie-breaking by assignment history;
multi-policy scheduling benchmark.}
& \cell{Rawat, Maitri\\(Developer)}
\\
\hline
\cell{\textbf{Phase II:} core orchestrator pipeline; Express REST API; job execution over the Docker
Engine socket; crash recovery.\\[2pt]
\textbf{Phase III:} ML-assisted autoscaling (linear-regression load forecaster); per-job cancellation of
running containers; automatic retries; first scheduling-comparison feature.}
& \cell{Joshi, Nikhil\\(Developer)}
\\
\hline
\cell{\textbf{Phase II:} React 19 operator UI; seeded workload generator v2; project documentation.\\[2pt]
\textbf{Phase III:} autoscaling and comparison dashboard pages (charts, results tables); workload
generator refinements; README and user documentation.}
& \cell{Kunwar, Vansh\\(Developer, UI \& Workload Generator)}
\\
\hline
\end{tabularx}

\vspace{0.65cm}

% CHALLENGES
\sectiontitle{Challenges/Roadblocks (7 pts)}

\vspace{0.05cm}
\instruction{(Describe the challenges that you have faced or are facing so far
and how you plan to solve them. Max 300 words).}

\vspace{0.18cm}
\answerbox{4.05cm}{%
All of the following have been solved.
\begin{itemize}
    \item \textbf{Race conditions on worker capacity.} When many jobs start and finish at once, two
    transactions could read the same free capacity and both reserve it. We lock the job and worker rows
    inside one transaction, add database \texttt{CHECK} constraints as a second guard, and retry MySQL
    deadlock victims (up to 4 attempts) so callers never see the error.
    \item \textbf{Stack migration mid-project.} The first version used TypeScript and PostgreSQL. To match
    the course stack we moved to JavaScript and MySQL, rewriting the locking queries for InnoDB and
    re-running every test against MySQL.
    \item \textbf{Merging two autoscalers.} The ML autoscaler was first built as a separate module that
    deleted workers and competed with the reactive one, and a merge commit reverted earlier features.
    We restored the lost files and folded ML into the existing autoscaler as a strategy, with only one
    strategy active at a time.
    \item \textbf{Fair comparisons.} Docker timing varies by a few percent between identical runs. Both
    strategies (or all policies) run the same seeded workload on a reset cluster, results are only
    called comparable when the work checksums match, and gaps under 5\% count as ties. A backend restart
    mid-experiment once left orphaned jobs; on startup the backend now marks such runs as stopped and
    cleans them up.
    \item \textbf{Readable results.} Early result pages were hard to interpret, so the dashboard now
    states a plain-language verdict with the measured numbers behind it, and shows N/A instead of
    inventing values when data is missing.
\end{itemize}}

\newpage

% ============================================================
% OUTCOMES AND PROGRESS
% ============================================================
\vspace*{0.30cm}

% OUTCOMES
\sectiontitle{Project Outcome/Deliverables (2 pts)}

\vspace{0.05cm}
\instruction{(Describe what are the key outcomes / deliverables of the project. Max 200 words).}

\vspace{0.35cm}
\answerbox{3.35cm}{%
\begin{itemize}
    \item A working orchestrator: seeded workload generation, four scheduling policies, three placement
    strategies, transaction-safe reservation and real Docker execution capped at each reservation.
    \item Reactive and ML-assisted autoscaling of logical workers, with every scaling decision recorded.
    \item Two reproducible experiments: a scheduling-policy benchmark and a Reactive vs ML autoscaling
    comparison, each with charts, a results table, a verdict and JSON export.
    \item A React dashboard that drives every pipeline stage and shows live monitoring, in light and
    dark themes and on phone-sized screens.
    \item Automated tests (241 backend, 66 frontend), Docker Compose setup and project documentation.
    \item Measured findings for the research write-up. For example, on a 25-job burst, Reactive kept the
    average wait 26\% shorter (7.4\,s vs 10.0\,s), while ML-assisted needed no extra workers
    (0 vs 19 extra worker-seconds).
\end{itemize}}

\vspace{0.75cm}

% PROGRESS
\sectiontitle{Progress Overview (2 pts)}

\vspace{0.05cm}
\instruction{(Summarize how much of the project is done, what's behind schedule,
what's ahead of schedule. Max 200 words.)}

\vspace{0.25cm}
\answerbox{4.45cm}{%
\textbf{The project is 100\% complete.} Every planned phase is done, merged into the \texttt{main}
branch and tested: persistent job and worker management, workload generation, policy scheduling,
resource-aware placement, transactional reservation, Docker execution, monitoring and the dashboard.

\vspace{4pt}
\textbf{Ahead of schedule:} several features first planned as stretch goals were also delivered:
Round Robin preemption with requeue, per-job cancellation of running containers, automatic retries,
reactive autoscaling, the ML-assisted autoscaler, the Reactive vs ML comparison experiment, and a full
redesign of the dashboard for readability.

\vspace{4pt}
\textbf{Behind schedule:} nothing. The one item recorded for future work, recovery from the failure of a
whole worker, was outside the original scope.}

\newpage

% ============================================================
% CODEBASE, TESTING, DELIVERABLES
% ============================================================
\vspace*{0.20cm}

% CODEBASE
\sectiontitle{Codebase Information (2 pts)}

\vspace{0.05cm}
\instruction{(Repository link, branch, and information about important commits.)}

\vspace{0.20cm}
\answerbox{3.70cm}{%
\textbf{Repository:} \texttt{https://github.com/Pranjalgit1/OrchestrOS}\\
\textbf{Branch:} \texttt{main} (feature branches merged through pull requests); 52 commits.\\[3pt]
\textbf{Important commits:}
\begin{itemize}
    \item \texttt{29b56ac}, \texttt{c8ac87b} (21 Sep): persistent job/worker management; deterministic
    workload batches.
    \item \texttt{f51e899}, \texttt{725937c}, \texttt{58381f8} (21--22 Sep): policy scheduling with safe
    claiming; resource-aware placement; transactions that prevent over-allocation.
    \item \texttt{6e1a54c} (23 Sep): React dashboard becomes the main control interface.
    \item \texttt{64ad1f6} (5 Oct): migration from TypeScript/PostgreSQL to JavaScript/MySQL.
    \item \texttt{3161357}, \texttt{1dfa173}, \texttt{2ad0f34} (21 Oct): Round Robin preemption,
    cancellation and requeue; reactive autoscaling; MySQL deadlock retry.
    \item \texttt{13ce8e7}, \texttt{8b619f5} (23 Oct): ML-based autoscaling; scheduling comparison.
    \item \texttt{287f464} (28 Oct): ML-predictive strategy merged into the autoscaler.
    \item \texttt{2b2c86e} (29 Oct): Reactive vs ML autoscaling comparison (new module and migration).
    \item \texttt{ec92e7a} (11 Nov): dashboard redesign (31 files).
\end{itemize}}

\vspace{0.65cm}

% TESTING
\sectiontitle{Testing and Validation Status (2 pts)}

\vspace{0.05cm}
\instruction{(Provide information about any tests conducted)}

\vspace{0.30cm}
\noindent
\begin{tabularx}{\textwidth}{
|>{\raggedright\arraybackslash}p{0.27\textwidth}
|>{\raggedright\arraybackslash}p{0.16\textwidth}
|>{\raggedright\arraybackslash}X|
}
\hline
\textbf{Test Type} &
\textbf{Status (Pass/Fail)} &
\textbf{Notes}
\\
\hline
\cell{Backend unit tests} & \cell{Pass (184/184)} &
\cell{Scheduling order, placement, autoscaler decisions, comparison metrics, input validation.}
\\
\hline
\cell{Database integration tests (MySQL)} & \cell{Pass} &
\cell{Real transactions: concurrent cancellation, exclusive execution claims, capacity release,
requeue and \texttt{CHECK} constraints.}
\\
\hline
\cell{Docker integration tests} & \cell{Pass} &
\cell{Real containers capped at their reservation, Round Robin eviction, cancellation, kill-all and
memory exhaustion. All 241 backend tests pass, 0 skipped.}
\\
\hline
\cell{Frontend view and unit tests} & \cell{Pass (66/66)} &
\cell{Rendered views, chart geometry, number formatting, N/A for missing data.}
\\
\hline
\cell{Static check and production build} & \cell{Pass} &
\cell{Backend syntax check and Vite production build.}
\\
\hline
\cell{End-to-end system test} & \cell{Pass} &
\cell{25-job Reactive vs ML comparison and 2- and 4-policy benchmarks on a fresh database; work
checksums matched across runs.}
\\
\hline
\cell{UI and responsive check} & \cell{Pass} &
\cell{Headless browser at desktop (1440\,px) and phone (390\,px) widths, light and dark themes.}
\\
\hline
\end{tabularx}

\vspace{0.65cm}

% DELIVERABLES
\sectiontitle{Deliverables Progress (2 pts)}

\vspace{0.05cm}
\instruction{(Summarize the current status of all key project deliverables
mentioned earlier. Indicate whether each deliverable is completed,
in progress, or pending.)}

\vspace{0.25cm}
\answerbox{3.20cm}{%
\begin{itemize}
    \item Orchestration pipeline (generation, scheduling, placement, reservation, execution):
    \textbf{Completed}
    \item Reactive and ML-assisted autoscaling: \textbf{Completed}
    \item Scheduling-policy benchmark: \textbf{Completed}
    \item Reactive vs ML autoscaling comparison: \textbf{Completed}
    \item React dashboard and live monitoring: \textbf{Completed}
    \item Automated tests, Docker Compose setup and documentation: \textbf{Completed}
    \item Findings for the research write-up: \textbf{Completed} (measured results available and exportable)
\end{itemize}}

\end{document}
